% Første linje er dokumentklassen. Brug *altid* memoir!
\documentclass[a4paper,oneside,article]{memoir}

% Herunder indlæses forskellige pakker
\usepackage[utf8]{inputenc}  % Korrekt håndtering af æ, ø og å
\usepackage[T1]{fontenc}  % Korrekt håndtering af æ, ø og å
\usepackage{microtype}  % Typografisk magi! Giver bl.a. pænere orddeling
\usepackage{graphicx}  % Gør det muligt at indsætte billeder
\usepackage{amsmath}  % Giver adgang til uundværlige matematikting
\usepackage{mathtools}  % Retter ting i ams og giver ekstra muligheder
\usepackage[danish]{babel}  % Danske betegnelser og orddeling
\usepackage{lmodern}
\usepackage{alphabeta}
\usepackage{float} %Mere kontrol over billede-placering
\renewcommand{\danishhyphenmins}{22}  % Bedre dansk orddeling
\usepackage{siunitx} % Sientific notation for forskellige ting, med en regel om at den skal bruge cdot istedet for times. 
\sisetup{inter-unit-product = \cdot, 
        exponent-product=\cdot}\usepackage{derivative} % Til at skrive differentialligninger 
\usepackage{xcolor} % For at få flere tekst farver 
\usepackage[colorlinks=true]{hyperref} % For at få refrences man kan trykke på, slet [] hvis man vil have kasser i stedet for farvet tekst.
\usepackage{pdfpages} % Til at indsætte andre pdf'er
\usepackage{subfig} % For at have to billeder ved siden af hinanden
\usepackage{unicode-math}
\usepackage{pgffor} % For the loop


% Pænere toc indents 
\usepackage{tocloft}
\cftsetindents{section}{1em}{2em}
\cftsetindents{subsection}{4em}{0em}


% Indholdet i dit dokument starter her!
\begin{document}
%\author{Mikkel Moth Billing \\ 202307989}
%\title{Title}
\date{\today}
%\maketitle   % Indsætter overskriften
% Table of Contents indstillinger
\setcounter{tocdepth}{2}
\setcounter{secnumdepth}{1}
\hypersetup{linkcolor=black} % Laver tableofcontents sort.  
%\tableofcontents %Indsætter ToC
\hypersetup{linkcolor=red} % og så alle andre hyperrefs røde. 

% Til at sætte pdf'er ind i store dokumenter (kun pdf'er på mere end 1 sider)
\newcommand{\pdf}[1]{
\includepdf[pages=1,scale=0.8,pagecommand={\subsection{#1}  \label{asc:#1}}]{pdf/stat/#1.pdf}
\includepdf[pages=2-, scale=0.8]{pdf/stat/#1.pdf}
}

\pdf{S1}
\pdf{Problem A1}
\pdf{Problem A2}
\pdf{S2}
\pdf{Problem B1}
\pdf{S3}
\pdf{Problem C1}
\pdf{S4}
\pdf{Problem D1}
\pdf{Problem D2}
\pdf{Problem D3}
\pdf{S5}
\pdf{Problem E1}
\pdf{Problem E2}
\pdf{Problem E3}
\pdf{Problem E4}
\pdf{S6}
\pdf{Problem F1}
\pdf{Problem F2}
\pdf{Problem F3}
\pdf{Problem F4}
\pdf{Problem G1}
\pdf{Problem G2}
\pdf{Problem G3}
\pdf{Problem G4}
\pdf{Problem G5}
\pdf{S8 - Hoffmann Problems}
\pdf{S8}
\pdf{Problem H1}
\pdf{Problem H2}

\end{document}